
\vspace{5mm}
\begin{ledgroup}
\footnotesize
\pstart
\noindent
Die vorliegenden zwei Konzepte N.\;\ref{60112} und N.\;\ref{60220} widmen sich in ähnlicher Weise derselben statischen Fragestellung und sind mit sehr hoher Wahrscheinlicheit etwa gleichzeitig, in einer Zeit intensiver Beschäftigung mit dieser und verwandten Fragen der Statik (siehe u.\,a. N.\;\ref{intro_vanni-solutio-problematis} und N.\;\ref{60116}) entstanden.
\pend
\pstart
Bei N.\;\ref{60112} lassen die Ähnlichkeiten mit dem Schreibduktus und den Papiereigenschaften des Konzeptes N.\;\ref{60115} eine Entstehung im selben Zeitraum vermuten; Leibniz verfasste N.\;\ref{60115} mit Sicherheit zwischen Dezember 1689 und Oktober 1690, wahrscheinlich auf der Durchreise durch Oberitalien oder Österreich.
Ebenso ist das besonders weiße und feine Papier von N.\;\ref{60220} vermutlich italienischer oder österreichischer Fabrikation (das Wasserzeichen im Träger konnte bislang nicht identifiziert werden.)
Aus den genannten Gründen wird die Datierung von N.\;\ref{60115} (Dezember 1689 bis Oktober 1690) für N.\;\ref{60112} übernommen und aus den folgenden inhaltlichen Gründen auch auf N.\;\ref{60220} übertragen.
\pend
\pstart
Die auffallenden Symmetrien in Aufbau und Inhalt beider Stücke lassen auf eine etwa gleichzeitige Abfassung schließen.
Denn sowohl N.\;\ref{60112} als auch N.\;\ref{60220} beginnen mit derselben Annahme 
einer (an einer Stelle abgeflachten) Kugel, die entlang einer schiefen Ebene lediglich gleiten, aber nicht rollen kann. Es handelt sich um das komplementäre Gedankenexperiment zum Wiener Konzept N.\;\ref{60116} (30. April bis Mitte Mai 1690), worin Leibniz eine auf der schiefen Ebene rollende Kugel annahm.
Die zwei Konzepte weisen außerdem teils wörtliche Übereinstimmungen im jeweils ersten Absatz.
Ebenfalls beiden Texten gemein ist die Fragestellung, nämlich nach den Bedingungen des Gleichgewichts der Kugel, allerdings unter zwei verschiedenen Voraussetzungen: In N.\;\ref{60112} wird die Kugel von unten mittels eines Hebels gestützt, in N.\;\ref{60220} aber von oben, mittels eines Seils, an dem sie hängt.
In beiden Fällen führt schließlich derselbe Ansatz zur Lösung: Das Gleichgewicht wird analytisch gesucht, mithilfe der "Methodu\bracks{s} de Maximis et Minimis" (\refnote{FnC:35_10_14_005_Maximis} in N.\;\ref{60112}).
\pend
\end{ledgroup}
